\chapter{Conclusion and Future Scope}
\label{ch:conclusion}

\begin{chapterabstract}
This chapter draws the thesis together. It revisits the research question and
the evidence assembled in support of the proposed scheduler, examines the
architectural trade-offs and limitations honestly, and sketches the most
promising directions for extending semantic-aware compute allocation in edge
robotic vision-language systems.
\end{chapterabstract}

\section{Thesis Synthesis and Final Review}
\label{sec:concl-synthesis}

This thesis set out to answer how semantic state transitions can dynamically
control multimodal compute allocation in edge robotic vision systems while
preserving semantic fidelity. The answer it develops is a
\emph{monitor-then-allocate} architecture in which an always-on, lightweight
monitor maintains a composite semantic state from cheap detection, tracking,
and embedding signals; estimates a bounded, interpretable drift signal that
fuses embedding distance, object-set novelty, and tracking-identity churn
(Equation~\ref{eq:drift}); and drives a three-tier policy
(Equation~\ref{eq:policy}) that spends expensive vision-language reasoning
only when the scene has genuinely changed.

The central conceptual contribution is to name and target \emph{semantic
reasoning redundancy} --- the repeated re-derivation of a scene's meaning while
that meaning is unchanged --- as an efficiency axis orthogonal to the pixel-,
frame-, and parameter-level redundancies that prior work addresses. The
practical contributions are the fused-drift formulation, the temporal
semantic memory that immunises it against sensor noise, the interpretable
tiered policy, the modular edge-oriented reference implementation, the
Semantic Compute Efficiency metric that makes the fidelity--compute trade-off
measurable on a single scale, and the Semantic Observatory together with the
observe--hypothesise--modify--evaluate methodology by which scheduling policies
are discovered from evidence rather than guessed. The evaluation of
Chapter~\ref{ch:experiments} demonstrates the expected behaviour: the
scheduler retains the bulk of an every-frame oracle's fidelity while invoking
the VLM on a small fraction of frames, achieving the highest SCE among the
policies compared.

\section{Architectural Trade-offs and Computational Costs}
\label{sec:concl-tradeoffs}

The approach is not without cost, and three trade-offs deserve explicit
acknowledgement.

\paragraph{Decision overhead.} The monitor runs on every frame, so its cost is
paid unconditionally. The design is only worthwhile because that cost is far
below a VLM invocation; on hardware where detection, tracking, and embedding
are not comfortably real-time, the margin shrinks. The frame filter and
reduced-resolution inference of Section~\ref{sec:arch-hardware} exist
precisely to protect this margin.

\paragraph{Memory footprint, not memory savings.} Because the VLM remains
resident in accelerator memory to avoid reload latency, peak memory is
unchanged relative to an every-frame policy. The savings are in time and
energy. A deployment that is memory- rather than compute-bound would need a
different treatment, such as offloading the VLM between bursts.

\paragraph{Threshold sensitivity and silent drift.} The policy's behaviour
depends on $\tau_{\text{low}}$, $\tau_{\text{high}}$, and the weights
$(\alpha,\beta,\gamma)$. Set too conservatively, the system reasons too
often; too aggressively, it risks missing slow, sub-threshold
\emph{silent drift}. The accumulating monitor of
Section~\ref{sec:arch-drift} mitigates but does not eliminate this risk, and
the ablation of Section~\ref{sec:exp-ablation} shows the trade-off is real.

\section{Directions for Future Research}
\label{sec:concl-future}

Several extensions follow naturally from this work.

\begin{itemize}
  \item \textbf{Learned, adaptive thresholds.} The fixed thresholds could be
  replaced by a lightweight controller that adapts $\tau_{\text{low}}$ and
  $\tau_{\text{high}}$ online to the observed drift distribution or to a
  downstream task-error signal, removing manual tuning.
  \item \textbf{Task-conditioned drift.} The drift weights are currently
  global. Conditioning them on the robot's active task --- so that, for
  instance, the appearance of task-relevant object classes dominates the
  signal --- would make redundancy reduction goal-aware.
  \item \textbf{Predictive scheduling.} The monitor reacts to drift after it
  occurs. Forecasting imminent semantic change from motion and tracking
  trajectories could let the system pre-warm or pre-empt reasoning, smoothing
  latency spikes at transitions.
  \item \textbf{Closing the loop with control.} Integrating the scheduler into
  a full vision-language-action stack~\cite{Kim2024OpenVLA} and measuring its
  effect on end-task success --- not just perceptual retention --- would validate
  the approach where it ultimately matters.
  \item \textbf{Broader hardware and model coverage.} Because the scheduler
  treats the VLM as a black box, evaluating it across additional compact VLMs
  and embedded accelerators would establish the generality of the SCE gains.
  \item \textbf{Cross-domain real-world validation.} The invocation patterns
  anticipated in Section~\ref{sec:exp-crossdomain} --- across surveillance,
  traffic, number-plate recognition, crowd, sports (cricket, football), and
  live-streaming video --- should be measured at scale to confirm that a single
  default configuration generalises, or to derive per-domain drift-weight and
  threshold profiles where it does not.
  \item \textbf{Egocentric robotic deployment.} Robots perceive through a
  moving camera, so egomotion robustness is essential rather than optional.
  The camera-motion compensation of Section~\ref{sec:meth-egomotion}, which
  discounts embedding drift by the egomotion fraction $c_t$ and relies on the
  motion-robust object and tracking signals, is the key enabler here;
  validating and refining it on head-mounted and mobile-robot streams --- and
  integrating richer egomotion estimation (e.g.\ visual odometry) --- is a
  priority for real-world use.
\end{itemize}

In closing, this thesis argues that the path to deploying vision-language
reasoning on resource-constrained robots runs not only through smaller and
faster models, but through systems that understand \emph{when} reasoning is
worth its cost. Treating semantic stability as a first-class, measurable
signal --- and allocating compute in proportion to semantic novelty --- offers a
principled and model-agnostic route to that goal.
